\section{Related Literature}\label{sec:literature}

\subsection{Interdependence and chokepoints}

\citet{hirschman1945national} located influence in the losses a partner would suffer if trade stopped, and \citet{keohane1977power} separated sensitivity, the cost of a cutoff before anyone adjusts, from vulnerability, the cost that remains afterward. \citet{farrell2019weaponized} moved the argument to networks and attached two conditions. A state must have ``physical or legal jurisdiction over hub nodes'' (p.~56), and it must have the ``legal and regulatory institutions'' to exploit them (p.~57); where hubs are spread across jurisdictions, the states involved must coordinate (pp.~56, 58). \citet{farrell2023underground} trace how Washington built that capacity, and \citet{kalyanpur2019mobilizing} show that jurisdiction is itself a policy choice: harmonized EU rules let European firms escape the reach of US securities law. Semiconductors offer the clearest case. The foreign direct product rule let Washington force foreign chipmakers to choose between US tools and Huawei as a customer \citep{bown2020how}, and US firms routed sales through third countries until that rule closed the gap \citep[194--195]{chen2023wars}. Treating a toolmaker as a hub is our extension, since Farrell and Newman's evidence comes from finance and the internet. On that extension ASML meets the network condition, but jurisdiction over it is shared with Washington (\cref{sec:background}), and the European institutions that could turn it against the United States are collective (\cref{sec:council}).

These works disagree about whether control of a hub or a supply yields lasting power. For Farrell and Newman the power comes from network position and the institutions behind it, and targets' responses, such as isolating themselves or reshaping the network, come later \citep[76]{farrell2019weaponized}. \citet{gholz2021market} find that Japan absorbed China's 2010 rare-earth embargo through market adjustment, and argue that cost measures that ignore market structure overstate a sender's power. Work on chokepoints as depleting assets sides with Gholz and Hughes: activation induces indigenization and rerouting \citep{tan2025driven, rajan2026eroding}. Our model lets the exit premium change over time (\cref{sec:control}) and asks what follows. Following \citet{wagner1988economic}, we treat asymmetric dependence as a source of political influence only if it changes a party's credible alternative, and ask when Europe's alternative is worth anything. \citet{drezner1999sanctions} adds the relevant asymmetry among senders: allies rarely coerce each other, but when they do they win significant concessions, because they expect little future conflict.

\subsection{Alliance coercion and dependence}

Alliance theory already separates an ally's dependence from the credibility of its patron's threats. \citet{snyder1984security, snyder1997alliance} distinguishes the fear of abandonment from the fear of entrapment, and a protégé's need for its patron from the patron's own strategic interest in defending it. \citet{morrow1991alliances} treats asymmetric alliances as trades of autonomy for security, and \citet{lake2009hierarchy} treats the patron's authority as the price of protection already received, sustained partly by legitimacy. On those accounts compliance needs no fresh threat, a rival reading we take up in \cref{sec:predictions}. \citet{leeds2003alliance} finds that commitments invoked in war between 1816 and 1944 were honoured about three-quarters of the time, with major powers likelier to renege, and \citet{fang2014concede} show that a target that values its alliance can be led to concede by its ally's advice.

The coercion literature disagrees about which instrument worked in the canonical case. \citet{gerzhoy2015alliance} credits conditional threats of military abandonment, paired with assurance, for stopping West Germany's nuclear ambitions. \citet{lanoszka2018atomic} finds alliance coercion hard and argues that an ally's economic and technological reliance on Washington is the stronger instrument. Our model contains both kinds. \citet{pauly2024damned} shows why assurance is hard: what makes a threat credible undermines the promise not to punish compliance. \citet{blankenship2021price} finds that allies in locations Washington values contribute less to the common defence, consistent with less credible abandonment threats, and \citet[111--119]{blankenship2023burden} shows Iceland extracting compensation and fishing concessions for Keflavik. \citet{carson2025allies} argue that access to allied territory is a benefit of alliance that a less committed patron may have to buy with cash, and \citet{wolford2017alliances} show that a patron's threats against a politically valuable ally can lack credibility, which is our self-harming-sanctions channel.

Base politics studies a host that holds an asset its patron needs. Political change in hosts reopens basing deals \citep{cooley2008base}, compensation shapes base stability \citep{calder2007embattled}, and sovereignty transfers are incomplete contracts that stay open to renegotiation \citep{cooley2009contracting}. In 2009 Kyrgyzstan announced the closure of the US base at Manas alongside a Russian aid package, collected the first Russian payment, and then signed a new deal with Washington at a higher rent \citep[ch.~7]{cooley2012great}. The nonproliferation literature studies the reverse problem, a patron restraining an ally's sensitive exports. \citet{kroenig2009exporting} finds that sensitive nuclear transfers follow suppliers' strategic position and finds no support for commercial motives, while \citet{gray2012commercial} shows Bonn putting commercial rivalry first in its 1975 sale to Brazil, keeping the deal and promising only future caution. \citet{miller2014secret} shows that credible US sanctions deterred nuclear programmes among states dependent on Washington, so that successes stay off the record, and \citet{coe2015collusion} model rival superpowers colluding to coerce. We add the interaction between the patron's conditionality and a rival's bid, the requirement that conditionality attach to the specific issue, and conditionality as a choice with an assurance cost (\cref{sec:conditionality}).

\subsection{Rival bidders and geoeconomics}

That a contested client earns what the losing bidder would pay is old. It is the one-voter case of sequential vote buying \citep{groseclose1996buying}, extended to finite legislatures by \citet{banks2000buying}; \citet[Prop.~2]{dekel2008vote} show that the winner pays the loser's value net of the client's preference, a price pinned down only in undominated strategies. Menu auctions \citep{bernheim1986menu} and their application to lobbying \citep{grossman1994protection} give the simultaneous-offer analogue, and \citet{bueno2016competition} show that a dominant aid donor captures the surplus of an aid-for-policy deal while a rival donor's entry raises what it must pay. The outside-option principle supplies the threshold: a threat to walk away matters only when the alternative beats continued bargaining \citep{shaked1984involuntary, sutton1986noncooperative, binmore1989outside}. \citet{lacy2004theory} model sanctions as coercive linkage under two-sided incomplete information and show that threats likely to succeed work before sanctions are imposed. \citet{jehiel1996how} show that with identity-dependent externalities a seller extracts surplus even from buyers who do not acquire the good.

Recent geoeconomics locates power in threats. \citet{clayton2026framework} show that a hegemon gains by pooling threats across relationships, bounded by the target's participation constraint, and \citet{clayton2026putting} state the principle we build on: ``Threats generate power to the extent that they generate slack in the participation constraint.'' The term slack is theirs. Their pooled threats are enforcement threats that raise the target's inside option; threats to cut off a target that refuses appear only in a supplement (B.3.1), where they are partly self-defeating, and a two-hegemon extension (B.2) drives transfers to zero when the hegemons' threats overlap. In \citet{clayton2026theory} a hegemon gains by committing to limit its own coercion. \citet{broner2025hegemonic} and \citet{becko2025optimal} model alignment among many small countries facing committed instruments. The rival's effect is contested. Rival offers come when a protégé's relations with its patron sour \citep{yang2026major}, and a rival's entry raised the US price in the aid market \citep{bueno2016competition}, yet \citet{camboni2025spheres} find that China's growth can lead countries in the US sphere to offer Washington more.

Each of these works lacks a result we derive. In \citet{groseclose1996buying} and \citet{dekel2008vote}, whose bidding outcome is our benchmark (\cref{lem:bidding}), the client is passive; once the holder can propose and talks can rupture, the patron's dependence on the chokepoint returns through the holder's agenda power (\cref{prop:master}). No principal in a menu auction \citep{bernheim1986menu} can move the agent's alternative. Ours can, and every conditional threat it holds then draws on one budget equal to the holder's slack (\cref{prop:slack}). \citet{bueno2016competition} give rival donors carrots only. Once the patron has sticks and the rival has counter-threats, China's threats lower the patron's payoff and can leave the holder's loss unchanged, and China executes one on the equilibrium path only to crowd out other demands (\cref{prop:counter}). Linkage in \citet{lacy2004theory} works through preferences and information; in our complete-information game it works through the conditional part of each instrument, and a cut once executed is spent (\cref{lem:selflimit}). The threats in \citet{clayton2026framework} are pooled enforcement threats, never spent, and their contracts are non-exclusive. Ours are cut-off threats that must each be sequentially rational, the rival's bid sits inside the participation constraint, and the holder chooses one alignment. Where the hegemon of \citet{clayton2026theory} gains by binding itself, our patron chooses how conditional to make its guarantee and pushes conditionality past what holds the premium at zero only to buy concessions (\cref{prop:conditionality}), and it is the holder that may bind itself (\cref{rem:bind}).

\subsection{Commitment, control and prevention}

Anticipated shifts in power create commitment problems that can make bargaining fail under complete information \citep{fearon1995rationalist, powell2006war}, and \citet{krainin2017preventive} shows that slow but persistent shifts suffice and that the preventive act may be delayed. \citet{fearon1996bargaining}, an unpublished manuscript, studies bargaining over objects that shape future bargaining power and finds efficient, gradual transfers when power varies continuously with holdings, and breakdown when it jumps. Control of ASML is such an object. Because Fearon's logic carries weight in our analysis of control, we note that the manuscript has not been peer reviewed. Uncertainty about a rising state's investment can itself produce preventive war \citep{debs2014known}, an observer that monitors imperfectly waits until the shift looks imminent \citep{bas2016dynamic}, and potential power can instead be bought off with concessions \citep{spaniel2019bargaining}. \citet{benson2023commitment} show that allies can offset an alliance-induced shift with concessions, and that the conditions that make prevention attractive also make it unlikely to be carried out.

Control rights come from the theory of the firm: ownership allocates residual rights over an asset when contracts are incomplete \citep{grossman1986costs, hart1990property}. \citet{lake1996anarchy} orders relations from alliance to empire by these rights, \citet{frieden1994international} ties colonial control to investments whose rents are easy to seize, and in \citet{harstad2016market} conservation abroad is leased while conservation at home is bought, with leasing dominant when the seller's consumption value and the buyer's cost of protecting what it owns are large. Sold tools are a wasting lever, since a durable-goods seller cannot commit \citep{coase1972durability, bulow1982durable}, while control of maintenance works like leasing and restores commitment \citep{morita2004durable}; servicing is therefore a renewable lever (\cref{sec:holdup}). Leaders can tie their hands with audience costs and other public stakes \citep{schelling1960strategy, fearon1994domestic, fearon1997signaling}; in our model every threat and promise is staked in this way, under one commitment technology (\cref{sec:commitment}). The conditional part of the threat behind the patron's preventive demand for control is capped by the holder's current slack, so it comes while the holder is weak but expected to strengthen (\cref{prop:control,cor:windows}). The transfer destroys sovereignty value and invites China's punishment in order to avoid paying rents that are pure transfers, which makes it the analogue of preventive war, and it binds only as far as the control it confers can be enforced (\cref{cor:control}).
